تو یک توسعه‌دهندهٔ ارشد React، TypeScript و CSS و متخصص رابط‌های تعاملی، انیمیشن و طراحی Responsive هستی.

می‌خواهم یک بازی آموزشی کامل و حرفه‌ای برای نمایش الگوریتم Bubble Sort بسازی. خروجی باید از نظر ظاهر، چیدمان، منطق، انیمیشن و عملکرد دقیقاً مطابق مشخصات زیر باشد.

نام پروژه:
Bubble Sort Robot Game

==================================================
۱. قانون اصلی تصاویر
==================================================

من هفت تصویر کاراکتر به تو می‌دهم.

- تصویر اول متعلق به کاراکتر شمارهٔ 1 است.
- تصویر دوم متعلق به کاراکتر شمارهٔ 2 است.
- تا تصویر هفتم که متعلق به کاراکتر شمارهٔ 7 است.
- فقط از همین هفت تصویر استفاده کن.
- هیچ تصویر جدیدی تولید نکن.
- تصاویر را با هوش مصنوعی بازسازی یا ویرایش نکن.
- پس‌زمینه، رنگ، چهره و ظاهر تصاویر را تغییر نده.
- فقط اندازه و موقعیت آن‌ها را با CSS تنظیم کن.
- هر کاراکتر همیشه باید همراه عدد خودش حرکت کند.
- هنگام جابه‌جایی، تصویر، عدد، Glow و تمام اجزای وابسته به آن باید به‌صورت یک واحد حرکت کنند.
- تصاویر را در مسیرهایی مانند زیر قرار بده:

public/images/stage/robot-1.png
public/images/stage/robot-2.png
public/images/stage/robot-3.png
public/images/stage/robot-4.png
public/images/stage/robot-5.png
public/images/stage/robot-6.png
public/images/stage/robot-7.png

اگر نام یا ترتیب فایل‌های پیوست‌شده مشخص نبود، قبل از پیاده‌سازی فقط دربارهٔ نگاشت تصاویر ۱ تا ۷ سؤال کن.

==================================================
۲. فناوری پروژه
==================================================

- React
- TypeScript
- Vite
- CSS خالص و حرفه‌ای
- بدون کتابخانه‌های سنگین رابط کاربری
- برای انیمیشن‌های جابه‌جایی از Web Animations API یا CSS Animation استفاده کن.
- منطق الگوریتم باید از رابط کاربری جدا باشد.
- پروژه باید با دستورات زیر بدون خطا اجرا شود:

npm run build
npm run lint

اگر پروژه‌ای از قبل وجود دارد، همان پروژه را ادامه بده و آن را از صفر بازسازی نکن. اگر پروژه خالی است، یک پروژهٔ سبک React + TypeScript + Vite بساز.

==================================================
۳. آرایهٔ اولیه
==================================================

آرایهٔ اولیه بازی:

[4, 1, 7, 3, 6, 2, 5]

ترتیب کاراکترهای روی صحنه باید دقیقاً از همین آرایه پیروی کند.

هر عدد نمایندهٔ هویت یک کاراکتر است، نه فقط مقدار یک خانه. بنابراین هنگام Swap باید خود کاراکترها جابه‌جا شوند و عدد همیشه با تصویر خودش باقی بماند.

==================================================
۴. ظاهر کلی
==================================================

سبک طراحی:

- محیط تاریک و مدرن
- رنگ اصلی پس‌زمینه سرمه‌ای بسیار تیره
- نورهای Neon آبی، فیروزه‌ای، بنفش، صورتی، سبز و زرد
- پنل‌های شیشه‌ای تیره با Border روشن
- سایه‌های نرم و Glow کنترل‌شده
- ظاهر حرفه‌ای و مناسب ضبط در OBS
- رابط فارسی و RTL
- کد پایتون و متغیرهای برنامه LTR باشند.
- از طراحی شلوغ، متن‌های تزئینی و عنوان‌های غیرضروری پرهیز کن.

این بخش‌ها نباید وجود داشته باشند:

- Header بزرگ
- لوگوی بالای صفحه
- عنوان Bubble Sort Robot Game در بالای صفحه
- عبارت ROBOT LAB
- عبارت «ربات‌ها را مرتب کن»
- عبارت «هر ربات همیشه با عدد خودش حرکت می‌کند»
- وضعیت «آماده شروع» در Header
- دکمهٔ بزرگ Reel Mode در Header

فضای حذف‌شده باید به صحنهٔ بازی اختصاص داده شود.

==================================================
۵. صحنهٔ کاراکترها
==================================================

صحنه در بالای صفحه قرار بگیرد.

مشخصات:

- هفت کاراکتر در هفت ستون مساوی نمایش داده شوند.
- عدد هر کاراکتر در یک Badge نئونی بالای آن باشد.
- زیر هر کاراکتر، اندیس آن به‌شکل زیر نمایش داده شود:

i:0
i:1
i:2
...
i:6

- در پایین صحنه یک سکوی بیضی‌شکل نئونی قرار بده.
- کف پای تمام کاراکترها باید طبیعی روی سطح سکو قرار بگیرد.
- کاراکترها نباید روی هوا باشند.
- کاراکترها نباید بیش‌ازحد داخل سکو فرو بروند.
- اگر تصاویر فضای خالی داخلی دارند، با CSS و یک Foot Offset مناسب آن‌ها را پایین بیاور.
- برای هر تصویر در صورت نیاز Offset جداگانه تعریف کن، اما خود فایل تصویر را ویرایش نکن.
- Badge اندیس نباید روی کفش یا پای کاراکتر قرار بگیرد؛ آن را روی قسمت جلویی سکو نمایش بده.
- اندازهٔ کاراکترها می‌تواند براساس شماره متفاوت باشد، ولی همه باید داخل صحنه جا شوند.
- کاراکتر شمارهٔ ۷ می‌تواند بلندترین و شمارهٔ ۱ کوتاه‌ترین باشد.
- هیچ اسکرول داخلی داخل صحنه وجود نداشته باشد.

در حالت عمودی، بالای صحنه یک فضای خالی مناسب برای تنفس بصری وجود داشته باشد. این فضا بخشی از محیط بازی است و نباید با عنوان یا دکمه پر شود.

مقدار پیشنهادی فضای بالای Arena:

clamp(48px, 7dvh, 96px)

==================================================
۶. پنل کد پایتون
==================================================

دقیقاً این کد را نمایش بده:

def bubble_sort(a):
    n = len(a)
    for p in range(n - 1):
        swapped = False
        for i in range(n - 1 - p):
            if a[i] > a[i + 1]:
                a[i], a[i + 1] = a[i + 1], a[i]
                swapped = True
        if not swapped:
            break

شرایط:

- کد تمام‌عرض باشد.
- هر ۱۰ خط کد همیشه کامل دیده شوند.
- پنل کد اسکرول داخلی نداشته باشد.
- فونت Monospace و خوانا باشد.
- شمارهٔ خطوط ۱ تا ۱۰ نمایش داده شود.
- Syntax Highlight داشته باشد.
- خط جاری الگوریتم Highlight شود.
- هنگام Compare خط ۶ با رنگ زرد Highlight شود.
- هنگام Swap خط ۷ با رنگ سبز Highlight شود.
- Header پنل شبیه پنجرهٔ کدنویسی باشد.
- نام فایل bubble_sort.py نمایش داده شود.
- پنل کد در حالت موبایل و عمودی مستقیماً زیر صحنهٔ کاراکترها قرار بگیرد.

==================================================
۷. چیدمان عمودی و موبایل
==================================================

حالت موبایل را با این شرط فعال کن:

- orientation: portrait
- یا width <= 768px

در تمام حالت‌های موبایل و عمودی:

- هیچ محدودیت عرض براساس نسبت ۹:۱۶ ایجاد نکن.
- عرض برنامه نباید به calc(100dvh * 9 / 16) محدود شود.
- کل رابط باید 100% عرض صفحه یا Browser Source در OBS را بگیرد.
- max-width باید 100% باشد.
- صحنه و کد باید تمام‌عرض باشند.
- چیدمان کاملاً ستونی باشد.
- هیچ پنل کناری در حالت عمودی وجود نداشته باشد.

چیدمان نمای اول:

1. صحنهٔ کاراکترها
2. کد کامل پایتون

ارتفاع پیشنهادی:

- صحنه: 60dvh
- کد پایتون: حدود 40dvh منهای Padding و Gap

نمای اول باید دقیقاً شامل صحنه و کل کد پایتون باشد.

بعد از اولین اسکرول، این بخش‌ها به‌ترتیب نمایش داده شوند:

1. آرایهٔ فعلی
2. کنترل‌های بازی
3. تنظیمات صدا و موسیقی
4. مقایسهٔ فعلی
5. متغیرهای الگوریتم
6. پیشرفت Passها یا آمار بازی دستی

صفحه باید اسکرول عمودی داشته باشد، اما صحنه و کد نباید اسکرول داخلی داشته باشند.

مهم‌ترین اندازهٔ تست:

720 × 1280

در این اندازه بررسی کن:

- صحنه تمام عرض باشد.
- فضای خالی بالای صحنه حفظ شود.
- کف پا روی سکو باشد.
- کد زیر صحنه باشد.
- هر ۱۰ خط کد تا قبل از انتهای ارتفاع ۱۲۸۰ دیده شوند.
- پنل آرایه از زیر نمای اول شروع شود.
- اسکرول افقی وجود نداشته باشد.

همچنین این اندازه‌ها را بررسی کن:

- 360 × 640
- 390 × 693
- 800 × 1024
- 1080 × 1920

==================================================
۸. چیدمان دسکتاپ
==================================================

در دسکتاپ از فضای افقی استفاده کن.

پیشنهاد:

- صحنه و اطلاعات اصلی در بخش بزرگ‌تر
- کد و متغیرها در ستون کناری
- کنترل‌ها در پایین
- ظاهر دسکتاپ نباید به‌خاطر حالت موبایل خراب شود.
- تمام منطق و امکانات در دسکتاپ نیز فعال باشند.

==================================================
۹. پنل آرایهٔ فعلی
==================================================

آرایه را با هفت Chip نئونی نمایش بده:

[4, 1, 7, 3, 6, 2, 5]

- رنگ هر Chip با رنگ کاراکتر مربوط هماهنگ باشد.
- جفت در حال مقایسه Highlight شود.
- پنل آرایه در موبایل زیر پنل کد قرار بگیرد، نه بین کاراکترها و کد.
- ارتفاع این پنل فشرده باشد.
- پنل آرایه نباید باعث پایین رفتن کد پایتون در نمای اول شود.

==================================================
۱۰. اجرای خودکار Bubble Sort
==================================================

حالت نمایش خودکار باید واقعی و مرحله‌ای باشد.

رویدادهای الگوریتم:

- compare
- swap
- pass-complete
- complete

قواعد:

- ابتدا دو عنصر مجاور مقایسه شوند.
- اگر عنصر چپ بزرگ‌تر بود، انیمیشن Swap اجرا شود.
- آرایه فقط بعد از فرود کامل کاراکترها به‌روزرسانی شود.
- نباید قبل از پایان انیمیشن ترتیب DOM تغییر کند.
- پس از پایان هر Pass، وضعیت Pass به‌روزرسانی شود.
- اگر در یک Pass هیچ Swap انجام نشد، الگوریتم زودتر متوقف شود.

کنترل‌های حالت خودکار:

- شروع
- توقف
- ادامه
- مرحلهٔ بعد
- اجرای مجدد
- اعداد تصادفی
- سرعت 0.5x
- سرعت 1x
- سرعت 1.5x
- سرعت 2x

Pause باید انیمیشن فعال را واقعاً متوقف کند و Resume همان انیمیشن را ادامه دهد.

==================================================
۱۱. انیمیشن‌های Swap
==================================================

حداقل چهار مدل انیمیشن جابه‌جایی بساز و آن‌ها را به‌صورت چرخشی اجرا کن:

1. پرش قوسی ساده
2. پرش همراه چرخش ۳۶۰ درجه
3. پرش دو‌مرحله‌ای
4. Flip یا چرخش سه‌بعدی

شرایط:

- دو کاراکتر هم‌زمان حرکت کنند.
- مسیر حرکت آن‌ها با هم تداخل بد نداشته باشد.
- یکی از کاراکترها هنگام عبور Z-Index بیشتری داشته باشد.
- Shadow هنگام پرش کوچک‌تر و کم‌رنگ‌تر شود.
- پس از فرود، کاراکترها دقیقاً روی جایگاه جدید قرار بگیرند.
- هیچ Jump یا تغییر ناگهانی بعد از پایان انیمیشن دیده نشود.
- چهار انیمیشن به‌صورت Deterministic و چرخشی استفاده شوند.
- Respect prefers-reduced-motion را رعایت کن.

==================================================
۱۲. بازی دستی واقعی
==================================================

یک حالت بازی دستی مستقل ایجاد کن.

روش استفاده:

- کاربر روی یک کاراکتر کلیک کند.
- یا با کلیدهای 1 تا 7 کاراکتر موردنظر را انتخاب کند.
- با ArrowLeft کاراکتر را یک خانه به چپ ببرد.
- با ArrowRight کاراکتر را یک خانه به راست ببرد.
- فقط جابه‌جایی با خانهٔ مجاور انجام شود.
- حرکت در لبهٔ چپ یا راست مجاز نباشد.
- در موبایل دکمه‌های چپ و راست لمسی وجود داشته باشند.

برای ارزیابی حرکت از تعداد Inversionهای آرایه استفاده کن.

اگر حرکت تعداد Inversion را کم کرد:

- بازخورد مثبت نمایش بده.
- افکت Correct اجرا کن.
- شمارندهٔ حرکت درست افزایش یابد.

اگر حرکت ترتیب را بدتر کرد:

- بازخورد منفی نمایش بده.
- افکت Wrong اجرا کن.
- شمارندهٔ اشتباه افزایش یابد.

آمار بازی دستی:

- تعداد کل حرکت‌ها
- حرکت‌های درست
- حرکت‌های اشتباه
- کاراکتر انتخاب‌شده
- اندیس فعلی
- آخرین حرکت

وقتی آرایه مرتب شد، حالت موفقیت نمایش داده شود.

==================================================
۱۳. متغیرهای الگوریتم
==================================================

در حالت خودکار این متغیرها را نمایش بده:

- n
- p
- i
- i + 1
- a[i]
- a[i + 1]
- swapped

مقادیر باید هم‌زمان با اجرای واقعی الگوریتم تغییر کنند.

در حالت دستی به‌جای آن‌ها نمایش بده:

- selected
- index
- moves
- correct
- wrong
- array

==================================================
۱۴. مقایسه و Pass
==================================================

پنل مقایسه باید نشان دهد:

- کدام دو مقدار مقایسه می‌شوند.
- اندیس i
- اندیس i + 1
- آیا Swap لازم است یا نه.

نمونه:

7 > 3 ، جابه‌جایی

برای Passها شش نقطه یا مرحله نمایش بده:

Pass 1 / 6

- Passهای کامل‌شده سبز یا روشن شوند.
- Pass جاری Glow داشته باشد.
- Passهای بعدی خاموش باشند.

==================================================
۱۵. صدا
==================================================

صدا در شروع خاموش باشد.

صدا فقط بعد از تعامل کاربر فعال شود تا با محدودیت Autoplay مرورگر سازگار باشد.

از Web Audio API استفاده کن و سیستم صوتی را ماژولار بساز.

رویدادهای صوتی:

- prepare
- compare
- swap
- jump
- flip
- land
- passComplete
- correctAnswer
- wrongAnswer
- success

امکانات:

- روشن و خاموش کردن افکت‌ها
- کنترل Volume
- روشن و خاموش کردن موسیقی مستقل از افکت‌ها
- موسیقی پس‌زمینهٔ الکترونیک بسیار ملایم
- هنگام Pause موسیقی متوقف شود.
- هنگام Resume ادامه پیدا کند.
- هنگام Complete موسیقی متوقف و صدای موفقیت اجرا شود.

ساختار صدا باید طوری باشد که بعداً بتوان فایل‌های صوتی واقعی را بدون تغییر منطق بازی جایگزین کرد.

==================================================
۱۶. حالت موفقیت
==================================================

وقتی آرایه مرتب شد:

- انیمیشن Confetti اجرا شود.
- یک پنل موفقیت نمایش داده شود.
- عبارت «مرتب شد!» نمایش داده شود.
- آرایهٔ نهایی زیر نمایش داده شود:

[1, 2, 3, 4, 5, 6, 7]

- در حالت دستی تعداد حرکت‌ها نیز نمایش داده شود.
- دکمهٔ «دوباره بازی کن» وجود داشته باشد.

==================================================
۱۷. معماری پیشنهادی
==================================================

فایل‌ها را منطقی تفکیک کن، مثلاً:

src/
  App.tsx
  App.css
  bubbleSortEngine.ts
  audioEvents.ts
  components/
    RobotStage.tsx
    CodePanel.tsx

Bubble Sort Engine باید مستقل از React باشد و Stepهای الگوریتم را تولید کند.

RobotStage مسئول این موارد باشد:

- نمایش کاراکترها
- انتخاب کاراکتر
- انیمیشن Swap
- Pause و Resume انیمیشن‌ها
- موقعیت صحیح پاها
- Confetti

CodePanel مسئول این موارد باشد:

- نمایش کد
- Highlight خط فعال
- نمایش متغیرها

==================================================
۱۸. نکات مهم پیاده‌سازی
==================================================

- هیچ‌وقت برای جابه‌جایی، عدد را از تصویر جدا نکن.
- از index آرایه به‌عنوان React key کاراکتر استفاده نکن؛ از RobotId استفاده کن.
- هنگام Swap تا پایان انیمیشن ترتیب React State را عوض نکن.
- تغییر Mode هنگام اجرای انیمیشن نباید باعث خرابی شود.
- Reset و Random باید تمام انیمیشن‌های فعال را Cancel کنند.
- تایمرها و AudioContext در Cleanup آزاد شوند.
- هیچ اسکرول افقی ایجاد نشود.
- هیچ متن یا کنترل مهمی بریده نشود.
- تصاویر کشیده یا فشرده نشوند.
- از object-fit: contain استفاده کن.
- عرض حالت موبایل همیشه 100% باشد.
- هیچ محدودیت مصنوعی 9:16 روی عرض اعمال نکن.
- نسبت 9:16 فقط یک اندازهٔ هدف برای OBS است، نه محدودکنندهٔ عرض CSS.

==================================================
۱۹. تست نهایی
==================================================

قبل از تحویل موارد زیر را واقعاً آزمایش کن:

- Build
- TypeScript
- Lint
- شروع خودکار
- Pause وسط انیمیشن
- Resume
- Step
- Reset
- Random
- تغییر سرعت
- مرتب‌سازی کامل
- حالت دستی با کلیک
- حالت دستی با کیبورد
- حرکت اشتباه
- حرکت درست
- جلوگیری از خروج در لبه
- موفقیت حالت دستی
- تغییر Mode وسط اجرا
- روشن و خاموش کردن صدا
- Volume
- موسیقی
- چیدمان دسکتاپ
- 360×640
- 390×693
- 720×1280
- 800×1024
- 1080×1920

در 720×1280 حتماً با اندازه‌گیری واقعی DOM تأیید کن:

- Stage تمام عرض باشد.
- Code تمام عرض باشد.
- Code زیر Stage باشد.
- خط ۱۰ کد داخل Viewport باشد.
- Array بعد از Code شروع شود.
- پای کاراکترها روی سکو باشد.
- Badgeهای i روی کفش‌ها نیفتند.

==================================================
۲۰. نحوهٔ تحویل
==================================================

کار را فقط با توضیح یا نمونه‌کد تحویل نده؛ پروژه را کامل پیاده‌سازی کن.

پس از اتمام فقط یک گزارش کوتاه بده:

- فایل‌های تغییرکرده
- نتیجهٔ Build و Lint
- امکانات پیاده‌شده
- اندازه‌هایی که آزمایش شدند
- محدودیت احتمالی باقی‌مانده

تا زمانی که Build موفق نشده و چیدمان 720×1280 بررسی نشده، کار را تمام‌شده اعلام نکن.